\section{Class-Specific Decoder Framework}

FFollowing training of the general autoencoder, class-specific decoders were
introduced to investigate whether the shared latent representation could be
mapped differently according to diagnostic group. The purpose of this
framework was to separate the learning of a common representation from the
learning of class-dependent reconstruction mappings.

Let the trained general autoencoder consist of an encoder $E$ and general
decoder $D_G$. After training, the encoder was frozen and its parameters
were no longer updated. Separate decoders were then trained for the HC,
MCI, and AD diagnostic groups while retaining the same shared latent
representation $\mathbf{Z}\in\mathbb{R}^{T\times L}$:

\begin{equation}
    \mathbf{X}
    \xrightarrow{E\;(\mathrm{frozen})}
    \mathbf{Z}
    \longrightarrow
    \left\{
    \begin{array}{l}
        D_G \\
        D_{\mathrm{HC}} \\
        D_{\mathrm{MCI}} \\
        D_{\mathrm{AD}}
    \end{array}
    \right.
\end{equation}

For a subject belonging to diagnostic group $c$, the reconstruction
produced by the corresponding class-specific decoder can therefore be
written as
\begin{equation}
    \hat{\mathbf{X}}_c
    =
    D_c\left(E(\mathbf{X})\right),
    \qquad
    c\in\{\mathrm{HC},\mathrm{MCI},\mathrm{AD}\}.
\end{equation}
Only the parameters of the corresponding class-specific decoder were
updated during this stage of training.

Each class-specific decoder was trained using samples belonging to its
corresponding diagnostic group. The decoder architecture was kept identical
to that of the general autoencoder decoder so that differences between
models reflected learned parameter differences rather than differences in
network structure. The class-specific decoders were optimised using the
same reconstruction objective used for the corresponding general
autoencoder configuration.

Freezing the encoder is important for the interpretation of this framework.
Because all decoders receive latent representations generated by the same
mapping $E$, the latent coordinate system remains common across the general
and class-specific models. Differences between
$D_G$, $D_{\mathrm{HC}}$, $D_{\mathrm{MCI}}$, and $D_{\mathrm{AD}}$
can therefore be interpreted as differences in how a shared latent
representation is mapped back to regional BOLD activity, rather than as
differences caused by independently learned encoders or incompatible latent
spaces.

The framework was also applied to subjects with longitudinal observations.
For an earlier scan $\mathbf{X}_{t_1}$, the frozen encoder first produces
the latent representation
\begin{equation}
    \mathbf{Z}_{t_1}
    =
    E(\mathbf{X}_{t_1}).
\end{equation}
This representation can then be decoded using the decoder corresponding to
a later diagnostic state,
\begin{equation}
    \hat{\mathbf{X}}_{t_2}^{(c)}
    =
    D_c(\mathbf{Z}_{t_1}),
\end{equation}
where $c$ denotes the target diagnostic class observed at the later time
point. This enables transitions such as HC$\rightarrow$MCI and
MCI$\rightarrow$AD to be examined within the same latent space. The
resulting reconstructions are interpreted as target-class mappings of the
earlier latent representation rather than as direct predictions of future
disease progression.